\documentclass[11pt,a4paper]{article}
\usepackage[margin=0.75in]{geometry}
\usepackage{xcolor}
\usepackage{listings}
\usepackage{hyperref}
\usepackage{booktabs}

\lstset{
    basicstyle=\ttfamily\small,
    breaklines=true,
    keywordstyle=\color{blue},
    commentstyle=\color{gray},
    stringstyle=\color{red},
    backgroundColor=\color{gray!10},
}

\title{VisioLearn Backend: System Architecture for Frontend}
\author{Backend Team}
\date{\today}

\begin{document}
\maketitle

\section*{1. System Overview}

VisioLearn is an offline-first audio learning platform for visually impaired students in Uganda. The backend implements a \textbf{class-based multi-subject architecture} where:

\begin{itemize}
    \item \textbf{Classes} are managed by a single \textit{class teacher} who generates two shareable codes
    \item \textbf{Subject teachers} join classes using the teacher code and can teach \textbf{multiple subjects}
    \item \textbf{Students} join classes using the student code and access content from all subject teachers
    \item \textbf{Progress} is tracked per student, per subject, per class
\end{itemize}

The API provides 44 routes across 8 resource categories. All endpoints are JWT-authenticated and support role-based access control (RBAC).

\section*{2. User Roles \& Registration}

\subsection*{Four User Roles}
\begin{table}[h]
\centering
\small
\begin{tabular}{|l|p{3.5cm}|}
\hline
\textbf{Role} & \textbf{Capabilities} \\
\hline
\texttt{admin} & Full system access, debugging, user management \\
\hline
\texttt{class\_teacher} & Creates class, generates codes, views all progress \\
\hline
\texttt{subject\_teacher} & Uploads content for subjects, sees subject progress \\
\hline
\texttt{student} & Joins class, accesses content, logs progress \\
\hline
\end{tabular}
\end{table}

\subsection*{Three Registration Endpoints}

\begin{lstlisting}[language=bash]
POST /api/v1/auth/register/class-teacher
Body: {email, full_name, password, class_name}
Returns: user_id, class_id, student_code (SC-XXXX),
         teacher_code (TC-XXXX)

POST /api/v1/auth/register/subject-teacher
Body: {email, full_name, password, 
       teacher_code?, subject_name?}
Returns: user_id, subject_id, class_id (if joined)

POST /api/v1/auth/register/student
Body: {email, full_name, password, student_code}
Returns: user_id, class_id, class_name
\end{lstlisting}

\textbf{Key Point:} Subject teachers can teach multiple subjects in the same or different classes. There is \textbf{no single-subject restriction}.

\section*{3. API Endpoints}

\subsection*{Authentication (7 routes)}
\begin{itemize}
    \item \texttt{POST /login} — Returns \texttt{access\_token}, \texttt{refresh\_token}
    \item \texttt{POST /refresh} — Renew expired tokens
    \item Registration endpoints (see above)
\end{itemize}

All endpoints use Bearer token: \texttt{Authorization: Bearer <token>}

\subsection*{Content Management (6 routes)}
\begin{lstlisting}[language=bash]
POST /api/v1/notes/upload
Body: {title, subject_id, grade_level, 
       description?, duration_seconds?}
Returns: LessonNote with id, class_id, subject_id

GET /api/v1/notes/list
Returns: Array of accessible notes
         (filtered by role & class membership)

GET /api/v1/notes/{id}
Returns: Single note (with access checks)
\end{lstlisting}

\textbf{Content Scoping:} Each note is linked to \textbf{(class\_id, subject\_id, teacher\_id)}. Students see only notes from subjects in their enrolled classes.

\subsection*{Progress Tracking (6 routes)}
\begin{lstlisting}[language=bash]
POST /api/v1/progress
Body: {content_id, last_position_seconds, completed}
Stores: class_id, subject_id, teacher_id 
        (auto-extracted from note)

GET /api/v1/progress/me
Returns: Student's overall progress

GET /api/v1/progress/me/by-subject
Returns: {subject_id → [progress items]}

GET /api/v1/progress/by-subject/{subject_id}
Returns: All students' progress for that subject
         (subject_teacher only)
\end{lstlisting}

\subsection*{Class Management (5 routes)}
\begin{lstlisting}[language=bash]
GET /api/v1/classes/{id}
Returns: {class_name, class_teacher_id, 
          student_code, teacher_code}

GET /api/v1/classes/{id}/students
Returns: [{id, email, full_name, joined_at}]
         (class_teacher only)

GET /api/v1/classes/{id}/matrix
Returns: {
  students: [{id, full_name}],
  subjects: [{id, name}],
  completion: [[0-100%, ...], ...]
}
\end{lstlisting}

\section*{4. Data Model}

\textbf{Core Entities:}
\begin{itemize}
    \item \textbf{User} — email, role, hashed\_password
    \item \textbf{Class} — class\_name, class\_teacher\_id, student\_code, teacher\_code
    \item \textbf{ClassSubject} — class\_id, subject\_name, subject\_teacher\_id
    \item \textbf{ClassMembership} — class\_id, student\_id, joined\_at, left\_at (soft delete)
    \item \textbf{LessonNote} — title, description, duration\_seconds, class\_id, subject\_id, teacher\_id, status
    \item \textbf{ContentProgress} — student\_id, content\_id, class\_id, subject\_id, teacher\_id, last\_position\_seconds, completed, completion\_percentage
\end{itemize}

\section*{5. Key Design Decisions}

\subsection*{Auto-Generated Codes}
Class teachers receive two shareable codes:
\begin{itemize}
    \item \textbf{Student Code:} \texttt{SC-XXXX} format (261K combinations)
    \item \textbf{Teacher Code:} \texttt{TC-XXXX} format (261K combinations)
\end{itemize}
Subject teachers and students join by entering these codes (no admin overhead).

\subsection*{Multi-Subject Teaching}
One subject teacher can teach:
\begin{itemize}
    \item Multiple subjects in the same class
    \item The same subject across different classes
    \item A subject mix: Math, Science, English in Class 9A; Biology in Class 10B
\end{itemize}
This is enabled by having \textbf{no unique constraint} on \texttt{(class\_id, subject\_teacher\_id)}.

\subsection*{Content Scoping}
Every lesson note belongs to exactly one (class, subject, teacher) triple. When a student accesses content, the backend verifies:
\begin{enumerate}
    \item Student is a member of the class
    \item The class has the subject
    \item The subject is taught by the requesting user
\end{enumerate}

\subsection*{Progress Tracking}
Progress is logged with full context: \texttt{(student\_id, content\_id, class\_id, subject\_id, teacher\_id)}. This enables:
\begin{itemize}
    \item Class teachers to see a matrix: students vs. subjects
    \item Subject teachers to see only their subject's progress
    \item Students to see their own progress per subject
\end{itemize}

\section*{6. Frontend Integration Guide}

\subsection*{Authentication Flow}
\begin{enumerate}
    \item User selects role: class teacher, subject teacher, or student
    \item Call appropriate registration endpoint
    \item Store returned \texttt{access\_token}, \texttt{refresh\_token} in secure storage
    \item Include token in all subsequent requests: \texttt{Authorization: Bearer <token>}
\end{enumerate}

\subsection*{Class Teacher Flow}
\begin{enumerate}
    \item Register class → get student\_code and teacher\_code
    \item Display codes for students/subject teachers to join
    \item Periodically fetch \texttt{GET /classes/{id}/students} to list enrolled students
    \item Fetch \texttt{GET /classes/{id}/matrix} to show progress dashboard
\end{enumerate}

\subsection*{Subject Teacher Flow}
\begin{enumerate}
    \item Register using teacher\_code from class teacher
    \item Upload content via \texttt{POST /notes/upload} with subject\_id
    \item View subject progress via \texttt{GET /progress/by-subject/{subject\_id}}
\end{enumerate}

\subsection*{Student Flow}
\begin{enumerate}
    \item Register using student\_code from class teacher
    \item Fetch \texttt{GET /notes/list} to get available content
    \item Play audio, periodically call \texttt{POST /progress} to log position
    \item View \texttt{GET /progress/me/by-subject} for progress summary
\end{enumerate}

\section*{7. Password Requirements}

All password fields require:
\begin{itemize}
    \item Minimum 12 characters
    \item At least one uppercase letter
    \item At least one lowercase letter
    \item At least one digit
    \item At least one special character: \texttt{!@\#\$\%\&*()?:..}
\end{itemize}

\section*{8. Error Handling}

All errors follow REST conventions:
\begin{itemize}
    \item \textbf{400} — Invalid request (validation failed, malformed UUID, etc.)
    \item \textbf{401} — Unauthorized (missing or invalid token)
    \item \textbf{403} — Forbidden (role/permission check failed)
    \item \textbf{404} — Resource not found (class, subject, note doesn't exist)
    \item \textbf{422} — Unprocessable entity (schema validation failed)
    \item \textbf{500} — Server error (database failure, etc.)
\end{itemize}

Response format:
\begin{lstlisting}[language=json]
{"detail": "Error message or array of validation errors"}
\end{lstlisting}

\section*{9. Base URL \& Environment}

\textbf{Current Deployment:} \url{https://visiolearn-backend.onrender.com}

All endpoints use: \texttt{https://visiolearn-backend.onrender.com/api/v1/<endpoint>}

For local development: \texttt{http://127.0.0.1:8000/api/v1/<endpoint>}

\section*{Questions?}

Contact the backend team. All endpoints are live and tested. The API documentation is also available via Swagger at \texttt{/docs}.

\end{document}
